\documentclass{article}  

\usepackage{listings}
\usepackage{amsmath}  
\usepackage{amsfonts}  
\usepackage{amssymb}
\usepackage{tikz}
\usepackage{circuitikz}
\usepackage{venndiagram}
\usepackage{enumitem}
\usepackage{csquotes}
\usepackage{titlesec}
\usepackage{tabularx}
\usepackage{ltablex}
\usepackage{hyperref}
\usepackage{float}
\usepackage{xcolor}



\lstset{
  basicstyle=\ttfamily\small, % Monospaced font
  keywordstyle=\color{blue}\bfseries, % Keywords in blue and bold
  commentstyle=\color{green}, % Comments in green
  stringstyle=\color{red}, % Strings in red
  numbers=left, % Line numbers on the left
  numberstyle=\tiny\color{gray}, % Style for line numbers
  frame=single, % Box around the code
  breaklines=true, % Automatic line breaking
  tabsize=2, % Tab size
  numberbychapter=false,
}

\lstdefinestyle{flexhighlight}{
    language=bash,
    escapeinside={@}{@}, % Enables LaTeX code inside listings
    basicstyle=\ttfamily
}

\renewcommand{\lstlistingname}{}
\newcommand{\highlight}[1]{\textbf{\textcolor{red}{#1}}}


\begin{document}  

\title{ITF20006-1 25V Algoritmer og datastrukturer\\Oblig 4}
\author{Andreas Isaksen}  
\date{\today}  
\maketitle  
\pagebreak

\section*{Oppgave 3}
Utskrift fra oppgave 2:
\begin{lstlisting}[language=bash, caption={Utskrift fra oppgave 2}]
n? 1000000
max? 10000000
n               tA      tL      tL/tA
--------------------------------------------
1000000         99      217     2,2  
2000000         55      250     4,5  
3000000         51      199     3,9  
4000000         57      214     3,8  
5000000         55      226     4,1  
6000000         54      239     4,4  
7000000         59      299     5,1  
8000000         60      300     5,0  
9000000         80      407     5,1  
10000000        59      293     5,0  
--------------------------------------------
\end{lstlisting}
Sorteringen i oppgave 2 går vesentlig mye raskere i forhold til sorteringen til oppgave 1. Dette er grunnet \textit{Timsort} som brukes internt av \textit{Arrays.sort()} og \textit{Collections.sort}\\\\
Når vi har en delvis sortert liste, vil denne bruke betraktelig kortere tid på å sortere den. Dette er grunnet at \textit{Timsort} kan gjenkjenne forhåndssorterte sekvenser i en liste, og videre slår disse sammen.\\
Om vi har en totalt tilfeldig liste, vil ikke løsningen være like effektiv til å dele opp listen og sortere den, da den må gå gjennom hele listen rekursivt.\\
Forskjellene i effektiviteten på sorteringen er derfor veldig avhengige av listens struktur før sorteringen starter.


\end{document}